% Related Work -- owner: Yiyang Sun.  % Related Work -- owner: Yiyang Sun.  % Related Work -- owner: Yiyang Sun.  % Related Work -- owner: Yiyang Sun.  \input{related_work} from the main report;
% citations are in references.bib.

\section{Related Work}
\label{sec:related}

\paragraph{Lifelong MAPD and decoupled planning.}
Ma et al.~\cite{ma2017lifelong} formulate lifelong multi-agent pickup and
delivery (MAPD) and propose Token Passing (TP) and its task-swapping variant
TPTS. TP is a decoupled, prioritized scheme in the spirit of Cooperative
A*~\cite{silver2005cooperative}: agents are planned one at a time with
space-time A* against a shared record of the paths already committed by
others, and a committed path is never revised. Within this scheme the record
of other agents' paths acts only as a set of hard obstacles; the step cost
remains the distance. Other work trades this simplicity for more coordination.
Liu et al.~\cite{liu2019task} solve offline MAPD, with all tasks known in
advance, by assigning task sequences via a special traveling-salesman problem
before planning paths. Li et al.~\cite{li2021lifelong} propose Rolling-Horizon
Collision Resolution (RHCR), which repeatedly solves windowed MAPF problems and
replans all agents. Foresight stays inside TP: it keeps one-at-a-time planning
and fixed paths and changes only the A* step cost. The cost therefore remains
as cheap to compute as TP, and every difference between our methods (M0--M5)
comes from that one term.

\paragraph{Static traffic rules.}
A common way to reduce congestion is to impose a direction of travel offline.
Cohen et al.~\cite{cohen2015feasibility} add highways, sets of directed edges
that bounded-suboptimal MAPF solvers are encouraged to follow; our M1 adapts
this idea to TP by penalizing moves against a highway edge, a variant those
authors list as future work. Zhang et al.~\cite{zhang2024guidance} optimize the
edge weights of a ``guidance graph'' offline for a given map and show that the
resulting guidance improves the throughput of several lifelong MAPF
algorithms. Both are designed per map before execution and do not react to the
traffic that actually occurs. Foresight's directional term (M5) aims at the same
outcome, one-way flow in corridors, but derives it online from the committed
paths, so the preferred direction can differ between corridors and change over
time.

\paragraph{Learned and history-based guidance.}
Direction maps~\cite{jansen2008direction} learn from how agents have moved
through the world and bias later agents to move with the prevailing direction,
producing implicit cooperation. Like our heatmap baseline (M2), which penalizes
cells in proportion to a decayed count of past visits, this guidance looks
\emph{backward}: it summarizes traffic that has already happened. M2 exists to
test whether looking \emph{forward}, at where committed agents will be, is
better than looking backward at where agents have been.

\paragraph{Congestion-aware guidance from planned routes.}
The closest work is the traffic flow optimisation of Chen et
al.~\cite{chen2024traffic}. They observe that lifelong MAPF planners route
agents along free-flow shortest paths and so create congestion, and they
compute \emph{guide paths} that avoid it. Guide paths are planned one agent at a
time on a graph whose edge costs grow with the flow of guide paths already
planned: a vertex-congestion term, and a contraflow term
$c_e = f(v_1,v_2)\cdot f(v_2,v_1)$, the product of the numbers of guide paths
crossing edge $e$ in each direction. The guide paths are then refined
continually and followed softly by PIBT, which prefers neighbors close to the
guide path. Foresight shares the idea of turning other agents' plans into a
congestion cost, including a penalty on opposing traffic, but differs in four
ways.
\begin{enumerate}
  \item \textbf{Time.} Their flows count \emph{time-independent} routes: a
  guide path loads an edge regardless of when the agent will be on it. Our
  congestion is read from collision-free \emph{space-time} paths and is
  time-indexed: $\rho(v,t)$ counts reservations of $v$ only around the time the
  planned agent would arrive (M4), and $c_{\mathrm{opp}}(u\to v,t)$ counts only
  opposing moves close to the time of the move (M5). A corridor that is busy now
  but empty when the agent gets there is not penalized.
  \item \textbf{Where the cost acts.} Their congestion shapes a separate
  guidance layer that is refined every timestep and that the underlying
  solver may deviate from. Ours is the step cost of the same space-time A*
  that produces the final path; it is evaluated once, when an agent is
  planned, and no agent is ever replanned.
  \item \textbf{Source of the forecast.} Their flows come from guide paths,
  which are themselves an estimate of where agents will go. In TP the
  committed paths are fixed, so the reservations they produce are an exact
  record of where every agent will be until its path ends.
  \item \textbf{Task assignment.} They assume an external task assigner.
  We also study replacing TP's nearest-pickup rule with a congestion-aware
  score $S_{ij}$, and separate the contribution of routing from that of
  assignment in an ablation.
\end{enumerate}
Their method targets thousands of agents with PIBT, whereas we study the
setting of TP with tens of agents; we therefore do not compare against it
directly.

\paragraph{Task assignment in MAPD.}
TP assigns each requesting agent the open task whose pickup is nearest
\cite{ma2017lifelong}, and TPTS lets an agent take over a task if it can reach
the pickup sooner. Liu et al.~\cite{liu2019task} assign whole task sequences
offline. These rules use distance only. Our $S_{ij}$ score adds the forecast
congestion along the agent's route to the pickup and delivery, and the time a
task has been waiting, to keep tasks from starving.
 from the main report;
% citations are in references.bib.

\section{Related Work}
\label{sec:related}

\paragraph{Lifelong MAPD and decoupled planning.}
Ma et al.~\cite{ma2017lifelong} formulate lifelong multi-agent pickup and
delivery (MAPD) and propose Token Passing (TP) and its task-swapping variant
TPTS. TP is a decoupled, prioritized scheme in the spirit of Cooperative
A*~\cite{silver2005cooperative}: agents are planned one at a time with
space-time A* against a shared record of the paths already committed by
others, and a committed path is never revised. Within this scheme the record
of other agents' paths acts only as a set of hard obstacles; the step cost
remains the distance. Other work trades this simplicity for more coordination.
Liu et al.~\cite{liu2019task} solve offline MAPD, with all tasks known in
advance, by assigning task sequences via a special traveling-salesman problem
before planning paths. Li et al.~\cite{li2021lifelong} propose Rolling-Horizon
Collision Resolution (RHCR), which repeatedly solves windowed MAPF problems and
replans all agents. Foresight stays inside TP: it keeps one-at-a-time planning
and fixed paths and changes only the A* step cost. The cost therefore remains
as cheap to compute as TP, and every difference between our methods (M0--M5)
comes from that one term.

\paragraph{Static traffic rules.}
A common way to reduce congestion is to impose a direction of travel offline.
Cohen et al.~\cite{cohen2015feasibility} add highways, sets of directed edges
that bounded-suboptimal MAPF solvers are encouraged to follow; our M1 adapts
this idea to TP by penalizing moves against a highway edge, a variant those
authors list as future work. Zhang et al.~\cite{zhang2024guidance} optimize the
edge weights of a ``guidance graph'' offline for a given map and show that the
resulting guidance improves the throughput of several lifelong MAPF
algorithms. Both are designed per map before execution and do not react to the
traffic that actually occurs. Foresight's directional term (M5) aims at the same
outcome, one-way flow in corridors, but derives it online from the committed
paths, so the preferred direction can differ between corridors and change over
time.

\paragraph{Learned and history-based guidance.}
Direction maps~\cite{jansen2008direction} learn from how agents have moved
through the world and bias later agents to move with the prevailing direction,
producing implicit cooperation. Like our heatmap baseline (M2), which penalizes
cells in proportion to a decayed count of past visits, this guidance looks
\emph{backward}: it summarizes traffic that has already happened. M2 exists to
test whether looking \emph{forward}, at where committed agents will be, is
better than looking backward at where agents have been.

\paragraph{Congestion-aware guidance from planned routes.}
The closest work is the traffic flow optimisation of Chen et
al.~\cite{chen2024traffic}. They observe that lifelong MAPF planners route
agents along free-flow shortest paths and so create congestion, and they
compute \emph{guide paths} that avoid it. Guide paths are planned one agent at a
time on a graph whose edge costs grow with the flow of guide paths already
planned: a vertex-congestion term, and a contraflow term
$c_e = f(v_1,v_2)\cdot f(v_2,v_1)$, the product of the numbers of guide paths
crossing edge $e$ in each direction. The guide paths are then refined
continually and followed softly by PIBT, which prefers neighbors close to the
guide path. Foresight shares the idea of turning other agents' plans into a
congestion cost, including a penalty on opposing traffic, but differs in four
ways.
\begin{enumerate}
  \item \textbf{Time.} Their flows count \emph{time-independent} routes: a
  guide path loads an edge regardless of when the agent will be on it. Our
  congestion is read from collision-free \emph{space-time} paths and is
  time-indexed: $\rho(v,t)$ counts reservations of $v$ only around the time the
  planned agent would arrive (M4), and $c_{\mathrm{opp}}(u\to v,t)$ counts only
  opposing moves close to the time of the move (M5). A corridor that is busy now
  but empty when the agent gets there is not penalized.
  \item \textbf{Where the cost acts.} Their congestion shapes a separate
  guidance layer that is refined every timestep and that the underlying
  solver may deviate from. Ours is the step cost of the same space-time A*
  that produces the final path; it is evaluated once, when an agent is
  planned, and no agent is ever replanned.
  \item \textbf{Source of the forecast.} Their flows come from guide paths,
  which are themselves an estimate of where agents will go. In TP the
  committed paths are fixed, so the reservations they produce are an exact
  record of where every agent will be until its path ends.
  \item \textbf{Task assignment.} They assume an external task assigner.
  We also study replacing TP's nearest-pickup rule with a congestion-aware
  score $S_{ij}$, and separate the contribution of routing from that of
  assignment in an ablation.
\end{enumerate}
Their method targets thousands of agents with PIBT, whereas we study the
setting of TP with tens of agents; we therefore do not compare against it
directly.

\paragraph{Task assignment in MAPD.}
TP assigns each requesting agent the open task whose pickup is nearest
\cite{ma2017lifelong}, and TPTS lets an agent take over a task if it can reach
the pickup sooner. Liu et al.~\cite{liu2019task} assign whole task sequences
offline. These rules use distance only. Our $S_{ij}$ score adds the forecast
congestion along the agent's route to the pickup and delivery, and the time a
task has been waiting, to keep tasks from starving.
 from the main report;
% citations are in references.bib.

\section{Related Work}
\label{sec:related}

\paragraph{Lifelong MAPD and decoupled planning.}
Ma et al.~\cite{ma2017lifelong} formulate lifelong multi-agent pickup and
delivery (MAPD) and propose Token Passing (TP) and its task-swapping variant
TPTS. TP is a decoupled, prioritized scheme in the spirit of Cooperative
A*~\cite{silver2005cooperative}: agents are planned one at a time with
space-time A* against a shared record of the paths already committed by
others, and a committed path is never revised. Within this scheme the record
of other agents' paths acts only as a set of hard obstacles; the step cost
remains the distance. Other work trades this simplicity for more coordination.
Liu et al.~\cite{liu2019task} solve offline MAPD, with all tasks known in
advance, by assigning task sequences via a special traveling-salesman problem
before planning paths. Li et al.~\cite{li2021lifelong} propose Rolling-Horizon
Collision Resolution (RHCR), which repeatedly solves windowed MAPF problems and
replans all agents. Foresight stays inside TP: it keeps one-at-a-time planning
and fixed paths and changes only the A* step cost. The cost therefore remains
as cheap to compute as TP, and every difference between our methods (M0--M5)
comes from that one term.

\paragraph{Static traffic rules.}
A common way to reduce congestion is to impose a direction of travel offline.
Cohen et al.~\cite{cohen2015feasibility} add highways, sets of directed edges
that bounded-suboptimal MAPF solvers are encouraged to follow; our M1 adapts
this idea to TP by penalizing moves against a highway edge, a variant those
authors list as future work. Zhang et al.~\cite{zhang2024guidance} optimize the
edge weights of a ``guidance graph'' offline for a given map and show that the
resulting guidance improves the throughput of several lifelong MAPF
algorithms. Both are designed per map before execution and do not react to the
traffic that actually occurs. Foresight's directional term (M5) aims at the same
outcome, one-way flow in corridors, but derives it online from the committed
paths, so the preferred direction can differ between corridors and change over
time.

\paragraph{Learned and history-based guidance.}
Direction maps~\cite{jansen2008direction} learn from how agents have moved
through the world and bias later agents to move with the prevailing direction,
producing implicit cooperation. Like our heatmap baseline (M2), which penalizes
cells in proportion to a decayed count of past visits, this guidance looks
\emph{backward}: it summarizes traffic that has already happened. M2 exists to
test whether looking \emph{forward}, at where committed agents will be, is
better than looking backward at where agents have been.

\paragraph{Congestion-aware guidance from planned routes.}
The closest work is the traffic flow optimisation of Chen et
al.~\cite{chen2024traffic}. They observe that lifelong MAPF planners route
agents along free-flow shortest paths and so create congestion, and they
compute \emph{guide paths} that avoid it. Guide paths are planned one agent at a
time on a graph whose edge costs grow with the flow of guide paths already
planned: a vertex-congestion term, and a contraflow term
$c_e = f(v_1,v_2)\cdot f(v_2,v_1)$, the product of the numbers of guide paths
crossing edge $e$ in each direction. The guide paths are then refined
continually and followed softly by PIBT, which prefers neighbors close to the
guide path. Foresight shares the idea of turning other agents' plans into a
congestion cost, including a penalty on opposing traffic, but differs in four
ways.
\begin{enumerate}
  \item \textbf{Time.} Their flows count \emph{time-independent} routes: a
  guide path loads an edge regardless of when the agent will be on it. Our
  congestion is read from collision-free \emph{space-time} paths and is
  time-indexed: $\rho(v,t)$ counts reservations of $v$ only around the time the
  planned agent would arrive (M4), and $c_{\mathrm{opp}}(u\to v,t)$ counts only
  opposing moves close to the time of the move (M5). A corridor that is busy now
  but empty when the agent gets there is not penalized.
  \item \textbf{Where the cost acts.} Their congestion shapes a separate
  guidance layer that is refined every timestep and that the underlying
  solver may deviate from. Ours is the step cost of the same space-time A*
  that produces the final path; it is evaluated once, when an agent is
  planned, and no agent is ever replanned.
  \item \textbf{Source of the forecast.} Their flows come from guide paths,
  which are themselves an estimate of where agents will go. In TP the
  committed paths are fixed, so the reservations they produce are an exact
  record of where every agent will be until its path ends.
  \item \textbf{Task assignment.} They assume an external task assigner.
  We also study replacing TP's nearest-pickup rule with a congestion-aware
  score $S_{ij}$, and separate the contribution of routing from that of
  assignment in an ablation.
\end{enumerate}
Their method targets thousands of agents with PIBT, whereas we study the
setting of TP with tens of agents; we therefore do not compare against it
directly.

\paragraph{Task assignment in MAPD.}
TP assigns each requesting agent the open task whose pickup is nearest
\cite{ma2017lifelong}, and TPTS lets an agent take over a task if it can reach
the pickup sooner. Liu et al.~\cite{liu2019task} assign whole task sequences
offline. These rules use distance only. Our $S_{ij}$ score adds the forecast
congestion along the agent's route to the pickup and delivery, and the time a
task has been waiting, to keep tasks from starving.
 from the main report;
% citations are in references.bib.

\section{Related Work}
\label{sec:related}

\paragraph{Lifelong MAPD and decoupled planning.}
Ma et al.~\cite{ma2017lifelong} formulate lifelong multi-agent pickup and
delivery (MAPD) and propose Token Passing (TP) and its task-swapping variant
TPTS. TP is a decoupled, prioritized scheme in the spirit of Cooperative
A*~\cite{silver2005cooperative}: agents are planned one at a time with
space-time A* against a shared record of the paths already committed by
others, and a committed path is never revised. Within this scheme the record
of other agents' paths acts only as a set of hard obstacles; the step cost
remains the distance. Other work trades this simplicity for more coordination.
Liu et al.~\cite{liu2019task} solve offline MAPD, with all tasks known in
advance, by assigning task sequences via a special traveling-salesman problem
before planning paths. Li et al.~\cite{li2021lifelong} propose Rolling-Horizon
Collision Resolution (RHCR), which repeatedly solves windowed MAPF problems and
replans all agents. Foresight stays inside TP: it keeps one-at-a-time planning
and fixed paths and changes only the A* step cost. The cost therefore remains
as cheap to compute as TP, and every difference between our methods (M0--M5)
comes from that one term.

\paragraph{Static traffic rules.}
A common way to reduce congestion is to impose a direction of travel offline.
Cohen et al.~\cite{cohen2015feasibility} add highways, sets of directed edges
that bounded-suboptimal MAPF solvers are encouraged to follow; our M1 adapts
this idea to TP by penalizing moves against a highway edge, a variant those
authors list as future work. Zhang et al.~\cite{zhang2024guidance} optimize the
edge weights of a ``guidance graph'' offline for a given map and show that the
resulting guidance improves the throughput of several lifelong MAPF
algorithms. Both are designed per map before execution and do not react to the
traffic that actually occurs. Foresight's directional term (M5) aims at the same
outcome, one-way flow in corridors, but derives it online from the committed
paths, so the preferred direction can differ between corridors and change over
time.

\paragraph{Learned and history-based guidance.}
Direction maps~\cite{jansen2008direction} learn from how agents have moved
through the world and bias later agents to move with the prevailing direction,
producing implicit cooperation. Like our heatmap baseline (M2), which penalizes
cells in proportion to a decayed count of past visits, this guidance looks
\emph{backward}: it summarizes traffic that has already happened. M2 exists to
test whether looking \emph{forward}, at where committed agents will be, is
better than looking backward at where agents have been.

\paragraph{Congestion-aware guidance from planned routes.}
The closest work is the traffic flow optimisation of Chen et
al.~\cite{chen2024traffic}. They observe that lifelong MAPF planners route
agents along free-flow shortest paths and so create congestion, and they
compute \emph{guide paths} that avoid it. Guide paths are planned one agent at a
time on a graph whose edge costs grow with the flow of guide paths already
planned: a vertex-congestion term, and a contraflow term
$c_e = f(v_1,v_2)\cdot f(v_2,v_1)$, the product of the numbers of guide paths
crossing edge $e$ in each direction. The guide paths are then refined
continually and followed softly by PIBT, which prefers neighbors close to the
guide path. Foresight shares the idea of turning other agents' plans into a
congestion cost, including a penalty on opposing traffic, but differs in four
ways.
\begin{enumerate}
  \item \textbf{Time.} Their flows count \emph{time-independent} routes: a
  guide path loads an edge regardless of when the agent will be on it. Our
  congestion is read from collision-free \emph{space-time} paths and is
  time-indexed: $\rho(v,t)$ counts reservations of $v$ only around the time the
  planned agent would arrive (M4), and $c_{\mathrm{opp}}(u\to v,t)$ counts only
  opposing moves close to the time of the move (M5). A corridor that is busy now
  but empty when the agent gets there is not penalized.
  \item \textbf{Where the cost acts.} Their congestion shapes a separate
  guidance layer that is refined every timestep and that the underlying
  solver may deviate from. Ours is the step cost of the same space-time A*
  that produces the final path; it is evaluated once, when an agent is
  planned, and no agent is ever replanned.
  \item \textbf{Source of the forecast.} Their flows come from guide paths,
  which are themselves an estimate of where agents will go. In TP the
  committed paths are fixed, so the reservations they produce are an exact
  record of where every agent will be until its path ends.
  \item \textbf{Task assignment.} They assume an external task assigner.
  We also study replacing TP's nearest-pickup rule with a congestion-aware
  score $S_{ij}$, and separate the contribution of routing from that of
  assignment in an ablation.
\end{enumerate}
Their method targets thousands of agents with PIBT, whereas we study the
setting of TP with tens of agents; we therefore do not compare against it
directly.

\paragraph{Task assignment in MAPD.}
TP assigns each requesting agent the open task whose pickup is nearest
\cite{ma2017lifelong}, and TPTS lets an agent take over a task if it can reach
the pickup sooner. Liu et al.~\cite{liu2019task} assign whole task sequences
offline. These rules use distance only. Our $S_{ij}$ score adds the forecast
congestion along the agent's route to the pickup and delivery, and the time a
task has been waiting, to keep tasks from starving.
